\documentclass[11pt]{article}
\usepackage{amsmath,amssymb,amsthm}
\usepackage{mathtools}
\usepackage{geometry}
\geometry{margin=1in}

\title{Multiplicity-LWE and CSL-Gated Multiplicity Dynamics:\\
A Unified Operator-Theoretic and Cryptographic Framework}

\author{Ryan O. Van Gelder \\ \\ Citizen Gardens \\ \\ The Foundation of Multiplicity \\ \\}
\date{March 19th, 2026}

\begin{document}

\maketitle

\begin{abstract}
We formalize a unified framework where multiplicity theory, prime-indexed operator
dynamics, and post-quantum cryptography interact coherently. The Multiplicity
Universal Constant $\Lambda_m$ is split into structural and analytic layers,
prime-sector operators are identified across the DRMM, PGF, and $\Xi(t)$
specifications, and a concrete Toy A instance is shown to instantiate a
standard Learning With Errors (LWE) problem over $\mathbb{Z}_{257}^{20}$.
We then introduce a CSL\_Gate that enforces consent, spectral drift bounds,
and norm budgets, and we state a gated convergence theorem (Theorem~\ref{thm:gated})
that wraps the Banach contraction guarantees of the ungated map.
\end{abstract}

\newpage
\tableofcontents
\newpage
\section{Prime-Indexed Structure and Canonical $\alpha$ Convention}

\subsection{Prime set and projectors}

Let $P_N = \{p_1, \dots, p_N\}$ be a finite set of primes. We work in a separable
Hilbert space $H$ and assume throughout an orthogonal prime-sector resolution.

\begin{assumption}[Orthogonal Prime Resolution]\label{ass:orthogonal-resolution}
There exists a family of orthogonal projectors $\{P_p\}_{p \in P_N} \subset B(H)$
such that
\[
P_p P_q = \delta_{pq} P_p,\qquad \sum_{p \in P_N} P_p = I.
\]
Consent operators $\Sigma(t)$ are required to be diagonal in this basis.
\end{assumption}

\subsection{Canonical $\alpha$ convention}

We adopt the standard number-theoretic convention for prime-weighted sums.

\begin{definition}[Prime-Weighted Zeta]\label{def:prime-zeta}
Let
\[
Z(\alpha) := \sum_{p} p^{-\alpha},
\]
where the sum runs over all primes $p$.
\end{definition}

\begin{lemma}[Convergence of Prime-Weighted Sums]\label{lem:prime-zeta-convergence}
The series
\[
\sum_{p} p^{-\alpha}
\]
converges if and only if $\alpha > 1$.
\end{lemma}

\begin{proof}
This is classical: $Z(\alpha)$ is dominated by the Riemann zeta function
$\zeta(\alpha)$ for $\alpha > 1$, and diverges for $\alpha \le 1$.
\end{proof}

\noindent
\textbf{Convention.} In DRMM, any previous condition of the form $\alpha < -1$ is
to be rewritten as $\alpha' = -\alpha > 1$, and the prime-weighted series is
written explicitly with $p^{-\alpha'}$. In PGF and $\Xi(t)$ specifications we
always write $\alpha > 1$.

\section{Unified Operator Binding Across DRMM, PGF, and $\Xi(t)$}

\subsection{Structural and analytic layers}

We distinguish between:

\begin{itemize}
    \item A \emph{structural layer} of PGF M-objects (e.g.\ $M_{\mathrm{fun}}$,
    $M_{\mathrm{int}}$, $\widehat{M}_c$) that encode prime-indexed recursion
    and functorial structure but do not themselves carry gains.
    \item An \emph{analytic layer} where the Multiplicity Universal Constant
    $\Lambda_m$ appears as a scalar certifying contraction under global and local
    policies (e.g.\ $\Lambda_{\mathrm{glob}}$ and $\Lambda_{\mathrm{loc}}$).
\end{itemize}

The structural layer is type-pure and functorial; all contraction and gain
properties live in the analytic layer.

\subsection{Operator identification}

Let $\{P_p\}_{p\in P_N}$ satisfy Assumption~\ref{ass:orthogonal-resolution}.
Define for each prime $p$ and time $t$:

\begin{definition}[Unified Prime-Sector Operators]\label{def:Up}
For each $p \in P_N$ and $t$, set
\begin{equation}
U_p(t) \;\stackrel{\mathrm{def}}{=}\; R_p(t) \;\stackrel{\mathrm{def}}{=}\;
\Lambda_m(t)\, p^{\alpha_p}\, P_p,
\tag{$\star$}
\end{equation}
where:
\begin{itemize}
    \item $\alpha_p > 1$ are real exponents (typically $\alpha_p \equiv \alpha$).
    \item $\Lambda_m(t)$ is a real-valued multiplicity gain, bounded in $t$,
    chosen to satisfy contraction conditions below.
    \item $P_p$ are orthogonal projectors with $\|P_p\|=1$.
\end{itemize}
\end{definition}

Equation~($\star$) identifies:
\begin{itemize}
    \item $U_p(t)$ from the $\Xi(t)$ specification,
    \item $R_p(t)$ from DRMM (prime-indexed morphisms),
    \item and the PGF-side prime-sector gain operator $\Lambda_m(t) p^{\alpha_p} P_p$.
\end{itemize}

\subsection{The $\Xi(t)$ evolution operator}

We define $\Xi(t)$ as a prime-weighted sum of the $U_p(t)$.

\begin{definition}[$\Xi(t)$ Operator]\label{def:Xi}
Let $w_p(t)$ be scalar weights. Define
\[
\Xi(t) := \sum_{p \in P_N} w_p(t)\, U_p(t)
= \Lambda_m(t) \sum_{p \in P_N} w_p(t)\, p^{\alpha_p}\, P_p.
\]
\end{definition}

\begin{assumption}[Bounded Weights]\label{ass:bounded-weights}
There exists $W < \infty$ such that for all $p \in P_N$ and all times $t$,
\[
|w_p(t)| \le W.
\]
\end{assumption}

From Definition~\ref{def:Up} and Assumption~\ref{ass:bounded-weights} we obtain
a simple norm bound:
\[
\|\Xi(t)\| \le |\Lambda_m(t)| \sum_{p \in P_N} |w_p(t)| p^{\alpha_p}
\le W \sup_t |\Lambda_m(t)| \sum_{p \in P_N} p^{\alpha_p}.
\]

\section{Ungated Contraction and the Role of $\Lambda_m$}

\subsection{Global and local contraction policies}

Let $\|S\| := \max_{p \in P_N} p^{\alpha_p}$ under the orthogonal resolution.
We define a global policy parameter $\gamma \in (0,1)$ and set
\[
\Lambda_{\mathrm{glob}} := \frac{\gamma}{\|S\|}.
\]

For a given non-linear map $T: H \to H$ with Lipschitz constant $L_T$,
we can also define a local Rayleigh-type gain $\Lambda_{\mathrm{loc}}$ so that
the effective contraction constant
\[
c := \sup_t \|\Lambda_{m,\mathrm{op}}(t)\|\, L_T
\]
stays below a desired margin $\varepsilon$. In practice we choose
$\Lambda_m(t)$ so that $c < \varepsilon$ and $\|\Xi(t)\| \le 1 - \varepsilon$,
yielding a uniform contraction.

\subsection{Ungated Banach contraction}

Consider the recursion
\begin{equation}\label{eq:ungated-recursion}
X_{t+1} = \Xi(t)\, X_t + \Lambda_m(t)\, T(X_t),
\end{equation}
where $T$ is non-linear with Lipschitz constant $L_T$ and $\Xi(t)$ is as in
Definition~\ref{def:Xi}. Suppose:
\begin{itemize}
    \item $\sup_t \|\Xi(t)\| \le 1 - \varepsilon$ for some $\varepsilon \in (0,1)$,
    \item $L_T$ is such that $c := \sup_t \|\Lambda_{m,\mathrm{op}}(t)\| L_T
    < \varepsilon$.
\end{itemize}
Then the combined map $F_t(X) := \Xi(t)X + \Lambda_m(t) T(X)$ is a contraction
in expectation or under suitable bounding conditions, and a standard Banach
argument yields a unique fixed point $X_\infty$ and exponential convergence:
\[
\|X_{t+1} - X_\infty\| \le \gamma \|X_t - X_\infty\|
\]
for some $\gamma < 1$ determined by $1-\varepsilon$ and $c$.\cite{postquantum-intro}

\section{Toy A: Multiplicity-LWE over $\mathbb{Z}_{257}^{20}$}

\subsection{Toy A parameter set}

We now fix a concrete Toy A instance that realizes a standard LWE problem
inside the multiplicity state space.

\begin{itemize}
    \item Prime set: $P_N = \{2,3,5,\dots,71\}$, the first 20 primes.
    \item Scale set: $S = \{0,1,2,3\}$, four scales.
    \item Ambient dimension: $n = |P_N| \cdot |S| = 80$.
    \item Exponent: $\alpha_p \equiv \alpha = 2$.
    \item Modulus: $q = 257$, a prime.
    \item Norm bound: $\|S\| = \max_{p \in P_N} p^{\alpha} = 71^2 = 5041$.
    \item Contraction target: $\gamma = 0.9$.
    \item Global gain: $\Lambda_{\mathrm{glob}} = \gamma / \|S\|
    \approx 0.9 / 5041 \approx 1.786 \times 10^{-4}$.
    \item Non-linear map: rounding modulo $q$ is non-expansive, so
    $L_T \le 1$ for our choice of $T$ (see below).
\end{itemize}

Thus $c = \Lambda_m L_T \lesssim 1.786 \times 10^{-4} \ll \varepsilon = 0.1$,
ensuring contraction.

\subsection{Scale layout and state space}

We index coordinates of $X_t \in H \cong \mathbb{C}^{80}$ by $(p,s)$ with
$p \in P_N$ and $s \in \{0,1,2,3\}$. We interpret the scales as follows:

\begin{itemize}
    \item $s=0$ (secret slice): holds the LWE secret vector
    $\mathbf{s} \in \mathbb{Z}_q^{20}$, one coordinate per prime $p$.
    This slice is frozen: neither $\Xi(t)$ nor $T$ updates it.
    \item $s=1$ (public random slice): holds the LWE sample row
    $\mathbf{a}_t \in \mathbb{Z}_q^{20}$ at time $t$, one coordinate per $p$.
    \item $s=2$ (inner product slice): holds
    $b_t = \langle \mathbf{a}_t, \mathbf{s} \rangle \bmod q$, computed
    from $s=0$ and $s=1$ via the linear part.
    \item $s=3$ (ciphertext slice): holds
    $b_t + e_t \bmod q$, where $e_t$ is a small error generated by $T$;
    this is the observable LWE sample.
\end{itemize}

\subsection{LWE-style non-linearity}

We define $T: H \to H$ to act only on the $s=2$ slice, injecting LWE-style noise.
Concretely, let $\chi$ be a centered binomial distribution with small parameter
$k=3$, approximating a discrete Gaussian with variance $\sigma^2 \approx 1.22$.
Let $e_t \in \mathbb{Z}_q$ be drawn from $\chi$ independently per coordinate
and time. Then $T(X_t)$ adds $e_t$ to the $s=2$ slice and writes the result
into the $s=3$ slice, followed by reduction modulo $q$ into a centered range.

The Lipschitz constant $L_T$ of $T$ can be bounded by $1 + B$ where $B$ is
a bound on the noise norm, but for small $e_t$ and centered reduction we can
take $L_T \le 1$ as a conservative overestimate in the toy.\cite{pqc-lwe}

\subsection{One-step recursion as LWE sampling}

Let $\mathbf{s}$ be fixed in the $s=0$ slice, and suppose the $s=1$ slice at
time $t$ contains $\mathbf{a}_t$. One step of the multiplicity recursion
\eqref{eq:ungated-recursion} is designed to implement:

\[
X_{t+1} = \Xi(t) X_t + \Lambda_m T(X_t),
\]

so that in particular:

\begin{align*}
\text{s=1 slice:} &\quad \mathbf{a}_{t+1} \text{ updated by key schedule},\\
\text{s=2 slice:} &\quad b_t = \langle \mathbf{a}_t, \mathbf{s} \rangle \bmod q,\\
\text{s=3 slice:} &\quad b_t + e_t \bmod q.
\end{align*}

Thus at each time $t$ the $s=1$ and $s=3$ slices contain a standard LWE sample
$(\mathbf{a}_t, b_t + e_t)$ with secret $\mathbf{s}$ and small error $e_t$,
over the modulus $q=257$ and dimension $n=20$.\cite{pqc-lwe}

\subsection{Inversion problem}

The inversion problem is:

\begin{quote}
Given the $s=3$ slices (ciphertext slice after $T$ steps) and the $s=1$ slices
for all $t$ (public random samples), recover the $s=0$ slice (the secret
$\mathbf{s}$).
\end{quote}

This is exactly an instance of the standard LWE problem over $\mathbb{Z}_q^n$
with $q=257$, $n=20$, and centered binomial noise $\chi$.\cite{pqc-lwe}

\section{CSL-Gated Dynamics and Consent}

\subsection{Consent-sector commutation}

We now introduce a CSL-style gate that constrains the evolution by consent,
spectral drift, and norm budgets.

\begin{lemma}[Consent-Sector Commutation]\label{lem:commutator}
Let $\{P_p\}_{p\in P_N}$ satisfy Assumption~\ref{ass:orthogonal-resolution}.
Let $U_p(t) = \Lambda_m(t)\, p^{\alpha_p}\, P_p \in B(H)$ as in Definition~\ref{def:Up},
and let $\Sigma(t)$ be a bounded operator on $H$. Then
\[
[\Sigma(t), U_p(t)] = 0 \quad \text{for all $p$ and $t$}
\]
if and only if $\Sigma(t)$ lies in the commutant of $\{P_p\}$, i.e.
\[
\Sigma(t) \in \{P_p\}' := \{A \in B(H): [A, P_p] = 0\ \forall p\}.
\]
\end{lemma}

\begin{proof}
By definition,
\[
[\Sigma(t), U_p(t)] = \Sigma(t)\Lambda_m(t) p^{\alpha_p} P_p - \Lambda_m(t) p^{\alpha_p} P_p \Sigma(t)
= \Lambda_m(t) p^{\alpha_p} [\Sigma(t), P_p].
\]
Since $\Lambda_m(t) p^{\alpha_p}$ is a nonzero scalar, $[\Sigma(t), U_p(t)] = 0$
iff $[\Sigma(t), P_p]=0$ for all $p$. This is precisely the condition
$\Sigma(t) \in \{P_p\}'$.
\end{proof}

Under Assumption~\ref{ass:orthogonal-resolution}, the algebra generated by
$\{P_p\}$ is:

\[
\mathrm{Alg}\{P_p\} = \left\{\sum_p \sigma_p P_p : \sigma_p \in \mathbb{C}\right\}.
\]

Choosing
\[
\Sigma(t) = \sum_p \sigma_p(t) P_p,\quad \sigma_p(t)\in\{0,1\}
\]
(binary consent per prime sector) ensures $[\Sigma(t), U_p(t)]=0$ for all $p,t$.

\subsection{Masked evolution and preserved contraction}

The masked evolution is
\[
\Sigma(t)\, \Xi(t)\, \Sigma(t)
= \sum_p \sigma_p(t)^2\, w_p(t)\, U_p(t)
= \sum_p \sigma_p(t)\, w_p(t)\, U_p(t),
\]
since $\sigma_p(t)^2 = \sigma_p(t)$ for $\sigma_p(t)\in\{0,1\}$. Thus masking
simply zeroes out non-consented sectors.

Moreover,
\[
\sup_t \|\Sigma(t)\Xi(t)\Sigma(t)\|
\le \sup_t \|\Xi(t)\| \le 1 - \varepsilon_\star,
\]
so masking can only preserve or tighten the contraction bound.

Under Assumption~\ref{ass:bounded-weights}, we have per-sector contraction
margins:
\[
|w_p(t)| \|U_p(t)\| \le 1 - \varepsilon_p
\quad\Rightarrow\quad
\varepsilon_p \ge 1 - W \sup_t |\Lambda_m(t)| p^{\alpha_p}.
\]
We can then define static drift weights
\[
\kappa_p = \frac{1}{\varepsilon_p}
\le \frac{1}{1 - W \sup_t|\Lambda_m(t)| p^{\alpha_p}},
\]
computable at initialization given $W$ and a bound on $|\Lambda_m(t)|$.

\subsection{CSL\_Gate and rate-limiter}

We assume a CSL\_Gate that enforces:

\begin{enumerate}
    \item A \textbf{consent quorum} $\mathrm{tr}(\Sigma(t)) \ge \theta_\Sigma$;
    otherwise the update is skipped.
    \item A \textbf{spectral drift bound} based on
    \[
    \delta_{\mathrm{drift}}(t) = \sum_p \kappa_p |\Phi(p,t) - \Phi(p,t-1)|,
    \]
    where $\Phi(p,t)$ is a prime-sector spectral statistic, and we require
    $\delta_{\mathrm{drift}}(t) \le \delta_{\max}$.
    \item A \textbf{norm budget} $\|\Xi(t)\|\le 1 - \varepsilon$; if violated,
    a rollback projects $X$ back onto a reference spectral slice
    $X \leftarrow \sum_\alpha E_\alpha(t_0) X$.
\end{enumerate}

We also define a rate-limiter with parameters $(R,N)$: in any window of $N$ steps,
at most $R$ rollbacks are allowed.

Let $\gamma<1$ be the contraction factor from the ungated Banach theorem,
and let $\delta_0$ be the initial error norm. To reach target precision
$\delta_{\mathrm{target}}$ we need at least
\[
M > \frac{\log(\delta_{\mathrm{target}}/\delta_0)}{\log(\gamma)}
\]
\emph{passing} steps. In a window of $N$ steps with at most $R$ rollbacks, at
least $N-R$ passes occur, so the liveness condition is
\begin{equation}\label{eq:rate-limiter}
N - R > \frac{\log(\delta_{\mathrm{target}}/\delta_0)}{\log(\gamma)}.
\end{equation}

\subsection{Gated convergence theorem}

We distinguish the ungated pure dynamics theorem from a new gated version.

\begin{theorem}[Ungated Convergence]\label{thm:ungated}
Suppose the recursion \eqref{eq:ungated-recursion} satisfies:
\begin{itemize}
    \item $\sup_t \|\Xi(t)\| \le 1 - \varepsilon$ for some $\varepsilon \in (0,1)$,
    \item $c = \sup_t \|\Lambda_{m,\mathrm{op}}(t)\| L_T < \varepsilon$,
\end{itemize}
for a map $T$ with Lipschitz constant $L_T$. Then there exists a unique fixed
point $X_\infty$ and a contraction factor $\gamma \in (0,1)$ such that
\[
\|X_{t+1} - X_\infty\| \le \gamma \|X_t - X_\infty\|
\]
for all $t$.
\end{theorem}

\begin{theorem}[Gated Convergence]\label{thm:gated}
Let $X_t$ evolve under the CSL-gated recursion: at each time $t$ either

\begin{itemize}
    \item the gate blocks (consent, drift, or norm-budget failure) and applies
    at most one rollback, or
    \item the gate passes and applies the update \eqref{eq:ungated-recursion}
    with masked $\Sigma(t)\Xi(t)\Sigma(t)$.
\end{itemize}

Assume:
\begin{enumerate}
    \item The ungated recursion satisfies Theorem~\ref{thm:ungated} with
    contraction factor $\gamma < 1$.
    \item In any sliding window of length $N$, at most $R$ rollbacks occur, and
    $(R,N)$ satisfy the rate-limiter constraint \eqref{eq:rate-limiter}.
\end{enumerate}
Then the gated trajectory converges to the same fixed point $X_\infty$ as the
ungated system, and there exists $t^\star$ such that
\[
\|X_t - X_\infty\| < \delta_{\mathrm{target}} \quad \text{for all } t \ge t^\star.
\]
\end{theorem}

\begin{proof}[Sketch]
Between rollbacks, passing steps contract errors by at least $\gamma$. Given at
most $R$ rollbacks in any $N$ steps, we have at least $N-R$ passing steps. The
inequality \eqref{eq:rate-limiter} guarantees enough passing steps to reduce
the error from $\delta_0$ below $\delta_{\mathrm{target}}$, and the masked map
is no less contractive than the ungated map.
\end{proof}

\section{Conclusion}

We have:

\begin{itemize}
    \item Fixed the $\alpha$ convention and unified the prime-indexed operator
    structure across DRMM, PGF, and $\Xi(t)$ via Equation~($\star$).
    \item Shown that Toy A instantiates a standard LWE instance over
    $\mathbb{Z}_{257}^{20}$ within the multiplicity state space, with small
    centered binomial noise.
    \item Introduced a CSL\_Gate with consent, drift, and norm-budget checks,
    and proven that, under a rate-limiter constraint, the gated dynamics converge
    to the same fixed point as the ungated system.
\end{itemize}

This provides a mathematically precise foundation for viewing multiplicity
dynamics both as a physically interpretable operator system and as a carrier
for post-quantum cryptographic hardness assumptions such as LWE.

\bibliographystyle{plain}
\begin{thebibliography}{9}

\bibitem{postquantum-intro}
Crypto Quantique.
\newblock An introduction to post-quantum cryptography.
\newblock 2019.

\bibitem{pqc-lwe}
J.~Newman.
\newblock Learning With Errors Post Quantum Cryptography.
\newblock GitHub repository, 2023.

\end{thebibliography}

\end{document}

